% Section 5 (Discussion) of "Conglomerate Mergers under Multiplicative Growth".
% 15 Nov 2026 draft from the front-matter doc (D1-D15 as ruled, Section 5 cards of 15 Nov with Simon's comments
% folded in: two-planners paragraph dropped, structure sentences dropped, chosen-rate paragraph shortened).
% Kept paragraphs are the current collaboration.tex text with the agreed trims. Simon writes the final prose over this.
% Labels expected elsewhere: sec:Background, eq:pooling_improvement, sec:model-economy, sec:model-conglomerates,
% sec:model-decisions, sec:model-calibration, sec:results-scope/-owners/-alpha/-distribution, tab:owner, tab:decomposition,
% app:A6, app:environment, app:cost, app:growth, app:A17, app:alpha.
\section{Discussion}
\label{sec:Discussion}

A recurring theme in the search for determinants of conglomerate formation is diversification to mitigate risk. As specialization of the Smithian \parencite{smith1776inquiry} and Ricardian \parencite{ricardo1817principles} kind raises efficiency, and as production has grown more roundabout through capital intensification \parencite{bawerk1921kapital}, the downside of an adverse run of shocks looms larger; the floor of Section~\ref{sec:model-economy} is where such a run ends, with the firm's capital written down to a minimum and recapitalized. Mitigating that downside has also become simpler \parencite{felton1971conglomerate}, since both internal and external means of diversification are available \parencite{goldberg1973effect}, their relative attractiveness depending on context. The model shows that pooling raises the time-average growth rate of every member and lowers the frequency with which a member is recapitalized at the floor, and that it does so between firms that share no market, technology or input. The shocks it insures against can be macroeconomic or political and specific to regions or sectors: the trade literature has argued that cross-country investment reduces country-specific risk and pays most between less correlated markets \parencite{ramondo2010role, ramondo2013proximity}, and conglomeration has been read as diversification across business cycles \parencite{amit1988concept} and across product life cycles \parencite{audretsch1989determinants}.

The floor is the model's exit. A firm that reaches it has been selected out: what survives is a configuration recapitalized at minimum scale, with its record but not its capital, and the floor-hit rate is the rate at which that happens. Pooling lowers a member's rate to a fraction of a standalone's (Section~\ref{sec:results-owners}) without changing anything about the firm's technology or market. Conglomerates are therefore a way of changing which firms are selected, and the unit, not the firm, becomes the object on which selection works. This complements the evolutionary account of corporate scope \parencite{teece1994understanding}, which explains what gets combined; here members share no market, technology or input, and combining pays anyway. The result is a disequilibrium one in the strict sense: decentralized formation under selection and finite information settles at a scope well short of the equal-members reference, and the gap is a property of the process, not of the parameters.

The model's scope can be set beside what diversified firms do, if the comparison is made at the right unit: a member occupies one market, so the analogue of $K$ is the number of distinct industries a firm or group operates in, not its size. \textcite{montgomery1994corporate} reports a mean of 10.7--10.9 four-digit SIC industries for the 500 largest US public companies in 1985--92, two thirds of them in five or more; the typical listed firm is far narrower \parencite{lichtenberg1992industrial}; and segment data, which cap the count at ten, understate both \parencite{villalonga2004diversification}. Four-digit industries within one two-digit sector are correlated, so the number of independent exposures is smaller than those counts. The model gives 3.5--7.7 members at imposed sharing rates of a tenth to a half, and 14--18 at the rate the units choose for themselves. Read as a pooling-only model should be read, this is the right side of the data: at moderate sharing rates it falls short of the breadth of the largest firms, as a model with no relatedness, no synergy and no market power ought to; that it overshoots at full pooling says that something outside the model keeps observed sharing rates below what unanimous members would choose, a point taken up below. Neither is a test; the model was not calibrated to breadth. Nor is the scale an artefact of the fifty markets: raising $M$ from 25 to 100 moves scope by half a member (Appendix~\ref{app:environment}), and units of four to eight members are nowhere near a ceiling that $M$ sets. The one prediction that does not rest on the model's stationary environment is the ordering across families: scope falls with the volatility of the environment through attrition, not through consent, while the willingness to pool rises with it. What limits conglomerates is consent on the recruits available and, under heavy tails, how often members are lost --- both observable, and neither the coordination cost that the Coasean reading would put first.

When units set their own sharing rate they settle at about 0.9 in every family and under every cost schedule (Section~\ref{sec:results-alpha}), far above anything a business group or an alliance pools and close to what a single firm does with its divisions. The reason is that in the model the sharing rate has no price: the coordination cost $\Phi$ depends on the number of members and not on how much they share, and a member's shocks are the same whether it keeps its return or passes it on, so $\alpha$ is a pure insurance parameter and the insurance-only optimum is near full pooling. In a real firm $\alpha$ is also an incentive parameter --- the deeper the reallocation, the weaker the link between a division's result and its resources, which is the socialism of internal capital markets that \textcite{scharfstein2000dark} and \textcite{rajan2000cost} describe --- and that cost is what the model deliberately leaves out. The chosen rate is therefore the insurance-only optimum, and the distance between it and observed practice measures the incentive cost the model does not carry. What the number is compared with matters too. A conglomerate merger in the legal sense is $\alpha = 1$: a wholly owned division's cash flow belongs to the parent, and what is partial is the reinvestment headquarters returns to it; a business group or a venture portfolio pools far less. The endogenous block says that if growth and survival were all that mattered, firms offered the choice would integrate almost fully --- which is what a merger is --- and the fixed-rate cells are the looser forms. The ordering is the model's ordering: integrated conglomerates pool more than groups, groups more than alliances, and the funds that pool least across a portfolio \parencite{boyd2023ergodic} are those whose portfolio firms least need the insurance to be deep.

Our results bear on the diversification discount. Early findings that diversified firms trade at a discount \parencite{lang1994tobins, berger1995diversification} have been challenged on endogeneity \parencite{campa2002explaining} and measurement \parencite{villalonga2004diversification}, and \textcite{kuppuswamy2016does} show the discount reversing in 2007--09. The model supplies a mechanism consistent with all three. The gain to a member's owners has two parts: faster growth in ordinary periods, and avoided ruin in the periods when a standalone of the same size in the same market would have been recapitalized at the floor. The second part is half of the gain under normal shocks and four fifths under heavy tails, and it is invisible to any comparison of typical periods: the median gap between members and standalones is a third to three fifths of the mean, and at low sharing rates it is zero while the mean is clearly positive (Table~\ref{tab:owner}). An evaluation by typical outcomes undervalues the arrangement by exactly the part that matters most where tails are heavy, and the methods of the discount literature are such evaluations --- comparisons of valuation multiples across firms in a cross-section of ordinary years. A crisis is the period in which the insurance pays out, which is what the 2007--09 reversal looks like from inside the model. The model also says why discount and premium coexist: units differ in who is in them (Appendix~\ref{app:A6}), whether the volatile side consented to the pooling, and how long a record admitted the members; a cross-section of conglomerates is a cross-section of contracts. Whether investors price the insurance ex ante is the empirical question the model sharpens rather than settles.

A related point concerns agency costs. The internal capital markets literature has documented how divisional rent-seeking corrupts capital allocation \parencite{scharfstein2000dark}, how headquarters with limited power misallocate capital across divisions \parencite{rajan2000cost}, and how free cash flow enables empire-building by self-interested managers \parencite{jensen1986agency}. These are real frictions. The coordination cost $\Phi(K)$ of Table~\ref{tab:parameters} captures the real resource costs of managing a larger organization in the tradition of \textcite{coase19883} and \textcite{williamson1967hierarchical} --- communication, monitoring and bureaucratic overhead that rise with scope --- which are conceptually distinct from agency costs arising from misaligned incentives between managers and shareholders. The model deliberately abstracts from the latter. By showing that an intrinsic growth gain from pooling exists in the complete absence of agency problems, it establishes a benchmark: the benefit of variance reduction under compounding is present regardless of governance quality, and agency costs determine how much of it is realized in practice. The crisis evidence of \textcite{kuppuswamy2016does} is consistent with this reading, suggesting that survival pressure in a crisis aligns managers' and owners' interests in the reallocation that pooling requires.

\textcite{lewellen1971pure} offers a purely financial rationale for conglomeration, where firms merge to mitigate financing frictions and to provide internal capital markets \parencite{boguth2022dissecting}. Our results allow a sharper distinction between this coinsurance channel and the multiplicative dynamics mechanism. The coinsurance rationale operates through a static distributional effect: combining imperfectly correlated cash flows enhances debt capacity. The mechanism we identify operates through temporal compounding as formalized in Equation~\ref{eq:pooling_improvement}: variance reduction raises the geometric growth rate of equity over time, independent of leverage. The two channels are complementary, but the growth-rate channel persists when firms face no financing frictions at all. Pooling can likewise be used to reduce the volatile returns inherent in R\&D, or to share its costs \parencite{felton1971conglomerate}.

The size distribution that pooling produces is more concentrated by one measure and less by every other measure the paper employs, and the one by which it is more concentrated is the one that matters least. Within markets firms become more equal: the within-market Herfindahl index falls and the tail exponent rises in every family, and since a unit holds one firm per market there is no horizontal market power anywhere in the model. What rises is the share of capital held by the largest owning units --- moderately, under heavy tails, and up to a sharing rate of one half; beyond that it falls again under normal and Laplace shocks as the units' members grow more equal (Appendix~\ref{app:A17}). The owners who become large and persistent are the insurers: their persistence is lower ruin hazard, not power, and the alternative at $\alpha = 0$ is an economy in which two thirds to four fifths of firms are recapitalized at the floor within a hundred periods --- configurations destroyed and rebuilt without end, which an earlier version of this paper called mobility. Rank mobility in a floored multiplicative economy is noise: ranks are nearly independent at a lag of five hundred periods at every sharing rate; pooling lowers the noise, so ranks move less, and no position is protected by anything but lower variance; markets stay contestable because any firm may form a unit. The same point explains why the tail exponent, not the pooled share, is the right measure of what pooling does: at a sharing rate of two per cent members already hold half the economy's capital and the tail is unchanged, because the exponent responds to how much of each shock is shared, not to how many firms share. Low rates do little to the tail and much to the hazard --- which is what a venture fund's modest pooling across its portfolio buys \parencite{boyd2023ergodic} --- and the mechanism applies to any multiplicative process, not only to financial capital \parencite{pietronero2001explaining, sibani2020human}. The one concern the results do raise is exclusion: incumbents refuse the small firms who would gain most from admission. Its remedy, if one is wanted, is a contract with an admission price or member-specific sharing, not antitrust.

The coordination cost is where the model is least informed by data, and the result is that it does not need to be. At equal cost for a unit of forty members, the four schedules move scope by less than one member; the level moves it by about 0.6 members per doubling, and twenty times the calibrated level leaves scope where the equal-members reference has fallen to meet it (Appendix~\ref{app:cost}). The exponential schedule produces the largest units only when units set their own rate, because the schedules are equal at forty members and the most convex of them is the cheapest below that. Nothing in Section~\ref{sec:Results} turns on the shape of $\Phi$.

The model assumes independent growth across markets. A cross-market correlation of 0.3 leaves scope within 0.2 of a member and every other outcome within the replication range (Appendix~\ref{app:environment}); \textcite{peters2022ergodicity} show the pooling gain surviving any correlation short of one. What the model does not cover is correlation that rises in a crisis \parencite{forbes2002no}, which is also when the floor is most likely to bind: the case in which the insurance is most valuable is the case in which it is thinnest. That is a real limitation and the first extension worth making.

The mechanism we identify operates alongside other forces that govern organizational form: selection determines which firms become conglomerates in the first place, agency frictions can corrupt the reallocation that makes pooling effective, operational integration generates synergies the model does not capture, and institutions shape the cost of coordination across borders and ownership structures. Our contribution is to isolate one of these forces under explicit assumptions. The crisis evidence of \textcite{kuppuswamy2016does} is where the pooling channel is most visible; identifying it in ordinary times would take considerably more demanding empirical work than the literature provides.

We emphasize that the model speaks to market structure rather than welfare. Reduced concentration is often associated with consumer welfare in standard models, but that connection requires assumptions about pricing, demand and allocative efficiency that the framework does not contain; horizontal concentration is excluded by construction, and what pooling concentrates is the ownership of capital across markets, not supply within any one. A welfare analysis would have to weigh, on the other side, a lower hazard of ruin, capital configurations kept in place, and less volatility drag on every member's growth. We state what such an analysis would weigh and do not make it; it requires models that integrate the growth-rate mechanism with explicit market interactions and a stance on the welfare criterion.
